\section{The set of natural numbers}

Throughout this section, $\N=\{1,2,3,\ldots\}$. An induction proof has
two separate obligations: establish an initial statement, and establish the
implication from an arbitrary statement to the next one. In the second
obligation, the induction hypothesis is an assumption used inside that
implication; it is not a substitute for the initial case. When induction
begins at an integer $m$ instead of at $1$, its conclusion concerns precisely
the integers $n\geq m$.

We write $\sum_{k=r}^{s}u_k$ for the sum of the terms indexed by
$r,r+1,\ldots,s$; if $r>s$, the sum is empty and equals $0$.
For integers $d\ne0$ and $M$, saying that $d$ divides $M$ means that
$M=dq$ for some integer $q$. The proofs of divisibility below identify this
integer explicitly.

\exercise{1.1}{Sum of the first squares}

\goal For every $n\in\N$, establish
\[
 \sum_{k=1}^{n} k^2=\frac{n(n+1)(2n+1)}{6}.
\]

\idea Increasing the upper limit from $n$ to $n+1$ adds exactly the
new term $(n+1)^2$. We substitute the proposed formula for the old sum,
then factor the resulting polynomial into the expression required at $n+1$.

\step{1. Specify the proposition and check the initial case.}
Let $P_n$ be the displayed identity. At $n=1$, its left side is $1^2=1$,
and its right side is
\[
 \frac{1(1+1)(2\cdot1+1)}{6}
 =\frac{1\cdot2\cdot3}{6}=1.
\]
Thus $P_1$ is true.

\step{2. State the hypothesis and the next target.}
Fix an arbitrary $n\in\N$ and assume $P_n$, namely
\[
 \sum_{k=1}^{n}k^2=\frac{n(n+1)(2n+1)}{6}.
\]
We must prove
\[
 \sum_{k=1}^{n+1}k^2
 =\frac{(n+1)((n+1)+1)(2(n+1)+1)}{6}
 =\frac{(n+1)(n+2)(2n+3)}{6}.
\]

\step{3. Add the new square and carry out the factorization.}
The induction hypothesis gives
\begin{align*}
 \sum_{k=1}^{n+1} k^2
 &=\left(\sum_{k=1}^{n}k^2\right)+(n+1)^2\\
 &=\frac{n(n+1)(2n+1)}{6}+(n+1)^2\\
 &=\frac{n(n+1)(2n+1)+6(n+1)^2}{6}\\
 &=\frac{(n+1)\bigl[n(2n+1)+6(n+1)\bigr]}{6}.
\end{align*}
To identify the bracket, expand its two products:
\[
 n(2n+1)+6(n+1)
 =2n^2+n+6n+6
 =2n^2+7n+6.
\]
The needed factorization is verified by multiplying it back out:
\[
 (n+2)(2n+3)
 =2n^2+3n+4n+6
 =2n^2+7n+6.
\]
Consequently,
\[
 \sum_{k=1}^{n+1} k^2
 =\frac{(n+1)(n+2)(2n+3)}{6},
\]
which is exactly $P_{n+1}$.

\answer{The initial case and the implication $P_n\Rightarrow P_{n+1}$
have both been proved. Mathematical induction therefore establishes the
formula for every positive integer $n$.}

\exercise{1.2}{An arithmetic progression}

\goal For every $n\in\N$, prove
\[
 \sum_{k=1}^{n}(8k-5)=4n^2-n.
\]
The summands at $k=1$ and $k=2$ are $8-5=3$ and $16-5=11$, so this
notation describes the progression in the exercise.

\idea Append its next term, $8(n+1)-5$, and compare the result with
the proposed expression evaluated at $n+1$.

\step{1. Initial case.}
For $n=1$, the left side is $3$, and the right side is
$4\cdot1^2-1=4-1=3$. Thus the statement holds at $1$.

\step{2. Induction hypothesis and target.}
Let $n\geq1$ and assume
\[
 \sum_{k=1}^{n}(8k-5)=4n^2-n.
\]
The required statement at the next index is
\[
 \sum_{k=1}^{n+1}(8k-5)=4(n+1)^2-(n+1).
\]
Before manipulating the sum, expand this target:
\begin{align*}
 4(n+1)^2-(n+1)
 &=4(n^2+2n+1)-n-1\\
 &=4n^2+8n+4-n-1\\
 &=4n^2+7n+3.
\end{align*}

\step{3. Compute the new sum.}
The appended term is
\[
 8(n+1)-5=8n+8-5=8n+3.
\]
Using the induction hypothesis,
\begin{align*}
 \sum_{k=1}^{n+1}(8k-5)
 &=\sum_{k=1}^{n}(8k-5)+\bigl(8(n+1)-5\bigr)\\
 &=(4n^2-n)+(8n+3)\\
 &=4n^2+(-n+8n)+3\\
 &=4n^2+7n+3\\
 &=4(n+1)^2-(n+1).
\end{align*}
This proves the next statement from the assumed one.

\answer{By induction, $3+11+\cdots+(8n-5)=4n^2-n$ for all $n\in\N$.}

\exercise{1.3}{Sum of cubes and the square of a sum}

\goal Prove, for every $n\in\N$, that
\[
 \sum_{k=1}^{n}k^3=\left(\sum_{k=1}^{n}k\right)^2.
\]

\idea First record the formula for the sum inside the square. Then
the asserted identity becomes a formula whose right side can be factored
after the next cube is added. All square expansions used here follow from
ordinary multiplication; no binomial theorem is being assumed.

\step{1. Establish the needed formula for the first-power sum.}
Write $T_n=1+2+\cdots+n$. We claim
\[
 T_n=\frac{n(n+1)}{2}\qquad(n\in\N).
\]
For $n=1$, both sides are $1$. If
$T_n=n(n+1)/2$ for some $n\geq1$, then
\begin{align*}
 T_{n+1}
 &=T_n+(n+1)\\
 &=\frac{n(n+1)}{2}+\frac{2(n+1)}{2}\\
 &=\frac{(n+1)(n+2)}{2}.
\end{align*}
This is the same formula with $n$ replaced by $n+1$. Induction therefore
proves the claim for every positive integer $n$.

It follows that the target identity can be written equivalently as
\[
 \sum_{k=1}^{n} k^3
 =\left(\frac{n(n+1)}{2}\right)^2
 =\frac{n^2(n+1)^2}{4}.
\]

\step{2. Check the initial case for the cube identity.}
For $n=1$, the sum of cubes is $1^3=1$, while
\[
 \frac{1^2(1+1)^2}{4}=
 \frac{1\cdot4}{4}=1.
\]
Thus the proposed formula holds at $1$.

\step{3. State the induction hypothesis and next target.}
Fix $n\geq1$ and suppose
\[
 \sum_{k=1}^{n} k^3=\frac{n^2(n+1)^2}{4}.
\]
We must show
\[
 \sum_{k=1}^{n+1} k^3
 =\frac{(n+1)^2(n+2)^2}{4}.
\]

\step{4. Append the next cube and factor.}
Using the hypothesis and a common denominator,
\begin{align*}
 \sum_{k=1}^{n+1}k^3
 &=\sum_{k=1}^{n}k^3+(n+1)^3\\
 &=\frac{n^2(n+1)^2}{4}+(n+1)^3\\
 &=\frac{n^2(n+1)^2+4(n+1)^3}{4}\\
 &=\frac{(n+1)^2\bigl[n^2+4(n+1)\bigr]}{4}\\
 &=\frac{(n+1)^2(n^2+4n+4)}{4}.
\end{align*}
Multiplication gives
\[
 (n+2)^2=(n+2)(n+2)
 =n^2+2n+2n+4=n^2+4n+4.
\]
Substituting this identity yields
\[
 \sum_{k=1}^{n+1}k^3
 =\frac{(n+1)^2(n+2)^2}{4}
 =\left(\frac{(n+1)(n+2)}{2}\right)^2
 =T_{n+1}^2.
\]
This is precisely the assertion required at $n+1$.

\answer{Mathematical induction proves
$1^3+2^3+\cdots+n^3=(1+2+\cdots+n)^2$ for every positive integer $n$.}

\exercise{1.4}{The sum of consecutive odd numbers}

\goal In part (a), use the first four sums to propose a formula. In
part (b), prove the proposed formula for every positive integer $n$.

\step{(a) Compute the requested examples and formulate the claim.}
Let
\[
 O_n=1+3+\cdots+(2n-1)=\sum_{k=1}^{n}(2k-1).
\]
The index $k$ runs from $1$ through $n$, so there are $n$ summands.
The requested calculations are
\[
\begin{array}{c|l|c|c}
 n&\text{sum}&O_n&n^2\\ \hline
 1&1&1&1\\
 2&1+3&4&4\\
 3&1+3+5&9&9\\
 4&1+3+5+7&16&16
\end{array}
\]
These examples suggest $O_n=n^2$. Checking four values gives a candidate
formula, but it does not establish what happens at every larger integer.

\step{(b) Check the initial case.}
At $n=1$, $O_1=1=1^2$, so the formula starts correctly.

\step{(b) State the induction hypothesis and next target.}
Fix $n\geq1$ and assume $O_n=n^2$. The target is
\[
 O_{n+1}=(n+1)^2.
\]
The new last summand is the $(n+1)$st odd number:
\[
 2(n+1)-1=2n+2-1=2n+1.
\]

\step{(b) Add the new odd number.}
The hypothesis gives
\begin{align*}
 O_{n+1}
 &=O_n+\bigl(2(n+1)-1\bigr)\\
 &=n^2+(2n+1)\\
 &=n^2+2n+1\\
 &=(n+1)^2.
\end{align*}
For the last equality, expanding $(n+1)(n+1)$ gives
$n^2+n+n+1=n^2+2n+1$.

\answer{The formula guessed in (a) is
$1+3+\cdots+(2n-1)=n^2$. Part (b) proves it for all $n\in\N$ by induction.}

\exercise{1.5}{A finite geometric sum}

\goal For every positive integer $n$, prove
\[
 1+\frac12+\frac14+\cdots+\frac1{2^n}
 =2-\frac1{2^n}.
\]
There are $n+1$ terms on the left, since the leading $1$ is $2^0$.
Equivalently, the sum is $G_n=\sum_{k=0}^{n}2^{-k}$.

\idea The next sum adds $1/2^{n+1}$. Expressing the old final
fraction with denominator $2^{n+1}$ leaves precisely the new final fraction.

\step{1. Initial case.}
For $n=1$, the left side is $1+1/2=3/2$, and the right side is
$2-1/2=4/2-1/2=3/2$. Thus the required statement holds at $1$.

\step{2. Hypothesis and target.}
Let $n\geq1$ and assume $G_n=2-1/2^n$. We must prove
\[
 G_{n+1}=2-\frac1{2^{n+1}}.
\]

\step{3. Add the next fraction with a common denominator.}
Because $2^{n+1}=2\cdot2^n$, we have
$1/2^n=2/2^{n+1}$. Hence
\begin{align*}
 G_{n+1}
 &=G_n+\frac1{2^{n+1}}\\
 &=2-\frac1{2^n}+\frac1{2^{n+1}}\\
 &=2-\frac2{2^{n+1}}+\frac1{2^{n+1}}\\
 &=2+\frac{-2+1}{2^{n+1}}\\
 &=2-\frac1{2^{n+1}}.
\end{align*}
All denominators are nonzero because powers of $2$ are positive.

\answer{The identity holds for every $n\in\N$ by mathematical induction.}

\exercise{1.6}{A difference of powers divisible by seven}

\goal Show that, for every $n\in\N$, there is an integer $q$ such that
\[
 11^n-4^n=7q.
\]
Here both $11$ and $4$ are raised to the power $n$.

\idea Express the next difference of powers as $11$ times the old
difference plus a term with a visible factor of $7$.

\step{1. Initial case and its integer witness.}
For $n=1$,
\[
 11^1-4^1=11-4=7=7\cdot1.
\]
The witness is the integer $q=1$, so $7$ divides the difference at $1$.

\step{2. Induction hypothesis and next target.}
Fix an arbitrary $n\geq1$ and assume $7$ divides $11^n-4^n$.
By the definition of divisibility, there is some $q\in\Z$ satisfying
\[
 11^n-4^n=7q.
\]
We must produce an integer $q'$ such that
$11^{n+1}-4^{n+1}=7q'$.

\step{3. Add and subtract a useful intermediate term.}
Exponent rules give $11^{n+1}=11\cdot11^n$ and
$4^{n+1}=4\cdot4^n$. Therefore
\begin{align*}
 11^{n+1}-4^{n+1}
 &=11\cdot11^n-4\cdot4^n\\
 &=11\cdot11^n-11\cdot4^n+11\cdot4^n-4\cdot4^n\\
 &=11(11^n-4^n)+(11-4)4^n\\
 &=11(7q)+7\cdot4^n\\
 &=7(11q+4^n).
\end{align*}
The second line adds $-11\cdot4^n+11\cdot4^n=0$, so it leaves the
quantity unchanged. Define $q'=11q+4^n$. Since $q$ is an integer and
$4^n$ is an integer, $q'$ is an integer. The last displayed equality
is the required divisibility witness at $n+1$.

\answer{Induction proves that $7$ divides $11^n-4^n$ for every positive
integer $n$.}

\exercise{1.7}{An exponential expression divisible by thirty-six}

\goal Prove that $36$ divides $7^n-6n-1$ for every $n\in\N$.

\idea Multiplying the old expression by $7$ nearly produces the
next one. The remaining difference will be $36n$.

\step{1. Initial case.}
At $n=1$,
\[
 7^1-6\cdot1-1=7-6-1=0=36\cdot0.
\]
Thus the divisibility statement holds, with integer witness $0$.

\step{2. Hypothesis and next target.}
Let $n\geq1$ and assume the current statement. Then some $q\in\Z$
satisfies
\[
 7^n-6n-1=36q.
\]
Our task is to express $7^{n+1}-6(n+1)-1$ as $36$ times an integer.

\step{3. Relate the next expression to the old one.}
First expand the expression at $n+1$:
\[
 7^{n+1}-6(n+1)-1
 =7\cdot7^n-6n-6-1
 =7\cdot7^n-6n-7.
\]
On the other hand,
\begin{align*}
 7(7^n-6n-1)+36n
 &=7\cdot7^n-42n-7+36n\\
 &=7\cdot7^n+(-42+36)n-7\\
 &=7\cdot7^n-6n-7.
\end{align*}
These two computations establish the exact identity
\[
 7^{n+1}-6(n+1)-1=7(7^n-6n-1)+36n.
\]
Now insert the induction hypothesis:
\begin{align*}
 7^{n+1}-6(n+1)-1
 &=7(36q)+36n\\
 &=36(7q+n).
\end{align*}
The number $7q+n$ is an integer because $q,n\in\Z$. It is therefore
the needed witness for divisibility at $n+1$.

\answer{The base case and the induction step show that
$36\mid(7^n-6n-1)$ for every positive integer $n$.}

\exercise{1.8}{Induction with a later starting index}

\goal Prove the inequality in (a) for all integers $n\geq2$, and
the inequality in (b) for all integers $n\geq4$.

\idea Use induction beginning at the first integer in each stated
domain. In part (b), the inequality from (a) will supply the comparison
needed after multiplying the induction hypothesis by $n+1$.

\step{(a) Initial case for $n^2>n+1$.}
At $n=2$, the inequality is
\[
 2^2=4>3=2+1.
\]
This gives the required base at the starting index $2$.

\step{(a) Induction hypothesis and next target.}
Let $n$ be an arbitrary integer with $n\geq2$ and assume $n^2>n+1$.
We must show
\[
 (n+1)^2>(n+1)+1=n+2.
\]
Expand and then use the assumed strict inequality:
\begin{align*}
 (n+1)^2
 &=n^2+2n+1\\
 &>(n+1)+2n+1\\
 &=3n+2.
\end{align*}
The inequality in the second line comes from adding the same number
$2n+1$ to both sides of $n^2>n+1$. Also,
\[
 (3n+2)-(n+2)=2n\geq4>0,
\]
so $3n+2>n+2$. Combining the two strict inequalities yields
$(n+1)^2>n+2$, as required.

\answer{Induction starting at $2$ proves $n^2>n+1$ for every integer
$n\geq2$.}

\step{(b) Initial case for $n!>n^2$.}
Recall that $n!=1\cdot2\cdots n$. At $n=4$,
\[
 4!=1\cdot2\cdot3\cdot4=24>16=4^2.
\]
Thus the required base is true at $4$.

\step{(b) Induction hypothesis and next target.}
Let $n\geq4$ be an integer and assume $n!>n^2$. The next target is
\[
 (n+1)!>(n+1)^2.
\]
The factorial definition gives $(n+1)!=(n+1)n!$. Since $n+1>0$,
multiplying the induction hypothesis by $n+1$ preserves its direction:
\[
 (n+1)!=(n+1)n!>(n+1)n^2.
\]
Because $n\geq4$ also implies $n\geq2$, part (a) applies and gives
$n^2>n+1$. Multiplying this inequality by the positive number $n+1$
gives the second comparison
\[
 (n+1)n^2>(n+1)(n+1)=(n+1)^2.
\]
Putting the comparisons together,
\[
 (n+1)!>(n+1)n^2>(n+1)^2.
\]
This is the needed induction step.

\answer{Induction beginning at $4$ proves $n!>n^2$ for every integer
$n\geq4$. The different starting indices are essential: the conclusions
of the two induction arguments are restricted to their stated domains.}

\exercise{1.9}{Exactly when an exponential exceeds a square}

\goal Determine the integers $n$ for which $2^n>n^2$, then justify
the infinite portion of the answer by induction.

\idea The strict inequality fails at some small positive integers,
so an induction argument cannot begin at $1$ and continue through every
subsequent index. We check the exceptional small cases separately, begin
the induction at $5$, and also address zero and negative integers because
the question says ``integers.''

\step{(a) Exclude every negative integer.}
Suppose $n\leq-1$ is an integer, and put $m=-n$. Then $m$ is a positive
integer and $n^2=m^2\geq1$. By the definition of negative powers,
\[
 2^n=2^{-m}=\frac1{2^m}.
\]
The denominator $2^m$ is a product of $m\geq1$ factors equal to $2$,
so $2^m\geq2$. Since the denominator is positive, taking reciprocals gives
$0<1/2^m\leq1/2$. Hence
\[
 0<2^n\leq\frac12<1\leq n^2.
\]
In particular, $2^n>n^2$ is false for every negative integer $n$.

\step{(a) Compute zero and the first five positive cases.}
\[
\begin{array}{c|r|r|c}
 n&2^n&n^2&2^n>n^2\text{?}\\ \hline
 0&1&0&\text{yes}\\
 1&2&1&\text{yes}\\
 2&4&4&\text{no: equality}\\
 3&8&9&\text{no: }8<9\\
 4&16&16&\text{no: equality}\\
 5&32&25&\text{yes}
\end{array}
\]
These calculations suggest that all integers at least $5$ work. They also
show that $0$ and $1$ are additional solutions, while $2,3,4$ are not.

\step{(b) Set up the induction beginning at $5$.}
Let $P_n$ denote $2^n>n^2$, now with the domain $n\geq5$.
The initial case is the checked inequality
\[
 2^5=32>25=5^2.
\]
Fix an arbitrary integer $n\geq5$ and assume $2^n>n^2$.
We must prove
\[
 2^{n+1}>(n+1)^2.
\]
Multiplication of the hypothesis by the positive number $2$ gives
\[
 2^{n+1}=2\cdot2^n>2n^2.
\]
To complete the step, we compare $2n^2$ with the next square.

\step{(b) Verify the remaining polynomial inequality.}
Expand the difference:
\begin{align*}
 2n^2-(n+1)^2
 &=2n^2-(n^2+2n+1)\\
 &=n^2-2n-1\\
 &=n(n-2)-1.
\end{align*}
Since $n\geq5$, we have $n-2\geq3>0$. Multiplication by $n-2$
therefore preserves the inequality $n\geq5$, and multiplication by $5$
preserves $n-2\geq3$. Thus
\[
 n(n-2)\geq5(n-2)\geq5\cdot3=15.
\]
It follows that
\[
 2n^2-(n+1)^2=n(n-2)-1\geq15-1=14>0.
\]
Consequently $2n^2>(n+1)^2$, and hence
\[
 2^{n+1}>2n^2>(n+1)^2.
\]
This proves the implication $P_n\Rightarrow P_{n+1}$ for every integer
$n\geq5$. Induction now proves $P_n$ for every such integer.

\answer{Taken literally over $\Z$, the complete solution set is
\[
 \{0,1\}\ \cup\ \{n\in\Z:n\geq5\}.
\]
If ``integers'' is restricted to the positive integers in the surrounding
discussion, the answer is $n=1$ or $n\geq5$. Zero is a solution whenever
zero is included in the domain. The negative integers and $2,3,4$ have all
been excluded, so no cases remain undecided.}

\exercise{1.10}{A sum whose lower and upper limits both change}

\goal Show that for every positive integer $n$,
\[
 (2n+1)+(2n+3)+\cdots+(4n-1)=3n^2.
\]

\idea Write the sum using an index before attempting induction.
When $n$ increases, the first old term disappears and two new terms enter.
Accounting for all three changes is the central step.

\step{1. Identify the summation limits and check the initial case.}
Define
\[
 S_n=\sum_{k=n}^{2n-1}(2k+1).
\]
At the lower limit $k=n$, the term is $2n+1$; at the upper limit
$k=2n-1$, it is $2(2n-1)+1=4n-1$. Consecutive values of $k$ increase
the term by $2$. There are $(2n-1)-n+1=n$ terms, so this is exactly
the sum in the question.

For $n=1$, both limits equal $1$. Thus
\[
 S_1=2\cdot1+1=3=3\cdot1^2.
\]
The formula is true at its first index.

\step{2. State the hypothesis and identify the next sum.}
Fix $n\geq1$ and assume $S_n=3n^2$. We must prove
$S_{n+1}=3(n+1)^2$. Replacing $n$ by $n+1$ in \emph{both} limits gives
\[
 S_{n+1}
 =\sum_{k=n+1}^{2(n+1)-1}(2k+1)
 =\sum_{k=n+1}^{2n+1}(2k+1).
\]
The old sum used the indices $n,n+1,\ldots,2n-1$. The new sum uses
$n+1,\ldots,2n-1,2n,2n+1$. Therefore it removes the old index $n$
and adds the indices $2n$ and $2n+1$.

\step{3. Record the common portion, including the empty case.}
Let
\[
 R_n=\sum_{k=n+1}^{2n-1}(2k+1).
\]
For $n=1$, this is the empty sum $\sum_{k=2}^{1}(2k+1)=0$;
the following decomposition remains valid in that case. For every $n\geq1$,
\begin{align*}
 S_n&=(2n+1)+R_n,\\
 S_{n+1}&=R_n+\bigl(2(2n)+1\bigr)+\bigl(2(2n+1)+1\bigr)\\
 &=R_n+(4n+1)+(4n+3).
\end{align*}
Solving the first equation for $R_n$ and substituting it into the second
establishes
\[
 S_{n+1}=S_n-(2n+1)+(4n+1)+(4n+3).
\]
In particular, the change cannot be described merely by appending one
term to $S_n$, since its lower limit also changes.

\step{4. Substitute the hypothesis and simplify the new value.}
\begin{align*}
 S_{n+1}
 &=3n^2-(2n+1)+(4n+1)+(4n+3)\\
 &=3n^2-2n-1+4n+1+4n+3\\
 &=3n^2+(-2+4+4)n+(-1+1+3)\\
 &=3n^2+6n+3\\
 &=3(n^2+2n+1)\\
 &=3(n+1)^2.
\end{align*}
This is exactly the next assertion.

\answer{Induction proves the stated sum equals $3n^2$ for all $n\in\N$.}

\exercise{1.11}{A valid induction implication without a true base}

\goal For $P_n$ defined by ``$n^2+5n+1$ is an even integer,'' prove
the implication requested in (a), determine its actual truth values in
(b), and explain what this teaches about induction.

\idea The difference between two consecutive expressions is even.
That fact proves the requested implication but does not supply a single
true initial statement. A separate parity calculation determines whether
any initial statement could exist.

\step{(a) Assume $P_n$ and state what must follow.}
Fix any $n\in\N$ and suppose $P_n$ is true. An integer is even when
it has the form $2q$ with $q\in\Z$, so the assumption says that
\[
 n^2+5n+1=2q
\]
for some integer $q$. We must prove $(n+1)^2+5(n+1)+1$ is also twice
an integer.

\step{(a) Expand the next expression and give its witness.}
Distributing the products gives
\begin{align*}
 (n+1)^2+5(n+1)+1
 &=(n^2+2n+1)+(5n+5)+1\\
 &=n^2+7n+7\\
 &=(n^2+5n+1)+(2n+6)\\
 &=(n^2+5n+1)+2(n+3).
\end{align*}
Substitute the assumed representation:
\[
 (n+1)^2+5(n+1)+1
 =2q+2(n+3)=2(q+n+3).
\]
Because $q+n+3\in\Z$, this expression is even. We have proved
$P_n\Rightarrow P_{n+1}$ for every $n\in\N$.

\step{(b) Determine the parity for an arbitrary $n$.}
Rewrite the expression as
\[
 n^2+5n+1=n(n+5)+1.
\]
Every integer is either even or odd. Consider each case explicitly.

If $n$ is even, write $n=2r$ for some $r\in\Z$. Then
\[
 n(n+5)+1=2r(n+5)+1=2\bigl[r(n+5)\bigr]+1.
\]
The bracket is an integer, so this has the form $2\ell+1$ and is odd.

If $n$ is odd, write $n=2r+1$ for some $r\in\Z$. Now
\[
 n+5=2r+1+5=2r+6=2(r+3),
\]
and hence
\[
 n(n+5)+1=n\cdot2(r+3)+1=2\bigl[n(r+3)\bigr]+1.
\]
Again the bracket is an integer, so the expression is odd.

Both cases show that $n^2+5n+1$ is odd for every $n\in\N$.
An odd integer cannot also be even: equality $2\ell+1=2q$ would imply
$1=2(q-\ell)$, whereas no integer multiplied by $2$ equals $1$.
Thus $P_n$ is false for every positive integer $n$.

\step{(b) Explain the missing logical requirement.}
The proposed initial statement $P_1$ is false because
\[
 1^2+5\cdot1+1=1+5+1=7,
\]
and $7$ is odd. Starting induction at a larger positive integer would
not help, since every one of the proposed starting statements is false.

Part (a) is nevertheless a valid implication. It says that \emph{if}
one of these expressions were even, then the next would be even too.
Its antecedent $P_n$ is never true, so the implication has no true
antecedent at which to start a chain of conclusions. In logical terms,
the implication is \emph{vacuously true}; it does not assert that either
$P_n$ or $P_{n+1}$ is actually true.

\answer{There is no $n\in\N$ for which $P_n$ is true, despite the validity
of $P_n\Rightarrow P_{n+1}$ at every index. The moral is that induction
requires both a true base case and the induction implication. Proving only
the implication proves none of the assertions unconditionally.}

\exercise{1.12}{The binomial theorem from factorials and induction}

\goal With
\[
 \binom{n}{k}=\frac{n!}{k!(n-k)!},\qquad 0\leq k\leq n,
\]
verify the first three cases, prove the identity in (b), and then prove
\[
 (a+b)^n=\sum_{k=0}^{n}\binom{n}{k}a^{n-k}b^k
 \qquad(n\in\N,\ a,b\in\R).
\]

\idea Multiplication by $a+b$ produces two sums. After shifting
the index in the sum multiplied by $b$, the terms with the same powers
have adjacent coefficients. Part (b) combines those coefficients into
the ones required for exponent $n+1$.

\step{Conventions and endpoint coefficients.}
For a positive integer $r$, $r!=1\cdot2\cdots r$, and $0!=1$.
All factorials used in denominators are therefore positive. In particular,
for each positive integer $n$,
\[
 \binom{n}{0}=\frac{n!}{0!n!}
 =\frac{n!}{1\cdot n!}=1,
 \qquad
 \binom{n}{n}=\frac{n!}{n!0!}
 =\frac{n!}{n!\cdot1}=1.
\]
The endpoint coefficients will be treated separately from the interior
ones, so no factorial of a negative integer will be needed.

For the polynomial expressions here, a power with exponent $0$ is an
empty product and denotes the constant $1$. Thus the expression $x^0$
in a monomial has value $1$ even when the value substituted for $x$ is
$0$. This convention identifies the endpoint monomials as $a^n$ and
$b^n$. No step below divides by $a$ or by $b$, so zero values of either
variable are permitted.

\step{(a) Verify the theorem at $n=1$.}
Here $1!=1$ and $0!=1$, so
\[
 \binom10=\frac{1!}{0!1!}=1,
 \qquad
 \binom11=\frac{1!}{1!0!}=1.
\]
The proposed right side is
\[
 \binom10a^1b^0+\binom11a^0b^1
 =1\cdot a\cdot1+1\cdot1\cdot b=a+b,
\]
which equals $(a+b)^1$. The identity therefore holds for every pair of
real numbers $a,b$ at $n=1$.

\step{(a) Verify the theorem at $n=2$ by multiplication.}
Since $2!=2$, the coefficients are
\[
 \binom20=\frac{2}{1\cdot2}=1,
 \qquad
 \binom21=\frac{2}{1\cdot1}=2,
 \qquad
 \binom22=\frac{2}{2\cdot1}=1.
\]
The distributive law and commutativity of real multiplication give
\begin{align*}
 (a+b)^2
 &=(a+b)(a+b)\\
 &=a(a+b)+b(a+b)\\
 &=a^2+ab+ba+b^2\\
 &=a^2+ab+ab+b^2\\
 &=a^2+2ab+b^2.
\end{align*}
Using the coefficients just computed, this is
\[
 \binom20a^2b^0+\binom21a^1b^1+\binom22a^0b^2,
\]
as required. This calculation does not appeal to a binomial expansion;
it establishes this particular expansion directly.

\step{(a) Verify the theorem at $n=3$ by another multiplication.}
The relevant factorials are $3!=6$, $2!=2$, $1!=1$, and $0!=1$.
Consequently
\[
 \binom30=\frac{6}{1\cdot6}=1,
 \quad
 \binom31=\frac{6}{1\cdot2}=3,
 \quad
 \binom32=\frac{6}{2\cdot1}=3,
 \quad
 \binom33=\frac{6}{6\cdot1}=1.
\]
Multiply the identity already established for the square by $a+b$:
\begin{align*}
 (a+b)^3
 &=(a+b)(a^2+2ab+b^2)\\
 &=a(a^2+2ab+b^2)+b(a^2+2ab+b^2)\\
 &=a^3+2a^2b+ab^2+a^2b+2ab^2+b^3\\
 &=a^3+(2+1)a^2b+(1+2)ab^2+b^3\\
 &=a^3+3a^2b+3ab^2+b^3.
\end{align*}
The coefficients are exactly $\binom30,\binom31,\binom32,\binom33$,
so this is the asserted formula at $n=3$, for all real $a,b$.

\step{(b) Set the range of the indices and obtain a common denominator.}
Fix $n\in\N$ and an integer $k$ with $1\leq k\leq n$. We must prove
\[
 \binom nk+\binom n{k-1}=\binom{n+1}{k}.
\]
The bounds on $k$ give $k-1\geq0$ and $n-k\geq0$, so every
factorial in the two coefficients is defined. From the factorial
definition,
\[
 (n-k+1)!=(n-k+1)(n-k)!,
 \qquad
 k!=k(k-1)!.
\]
The numbers $k$ and $n-k+1$ are both positive. Therefore the following
rewritings use nonzero factors and are valid:
\begin{align*}
 \binom nk
 &=\frac{n!}{k!(n-k)!}
 =\frac{n!(n-k+1)}{k!(n-k+1)!},\\
 \binom n{k-1}
 &=\frac{n!}{(k-1)!(n-(k-1))!}
 =\frac{n!}{(k-1)!(n-k+1)!}
 =\frac{n!k}{k!(n-k+1)!}.
\end{align*}

\step{(b) Add the numerators and identify the next coefficient.}
Both fractions now have denominator $k!(n-k+1)!$, and hence
\begin{align*}
 \binom nk+\binom n{k-1}
 &=\frac{n!(n-k+1)+n!k}{k!(n-k+1)!}\\
 &=\frac{n!\bigl[(n-k+1)+k\bigr]}{k!(n-k+1)!}\\
 &=\frac{n!(n+1)}{k!(n-k+1)!}\\
 &=\frac{(n+1)!}{k!((n+1)-k)!}\\
 &=\binom{n+1}{k}.
\end{align*}
The penultimate line uses $(n+1)!=(n+1)n!$ and
$n-k+1=(n+1)-k$. This proves the claimed identity for every integer
$k$ in the stipulated range $1,\ldots,n$. It is often called
Pascal's identity. This derivation uses only the definition of the
coefficients, not the binomial theorem being proved.

\step{(c) State the induction proposition and verify its base.}
Let $P_n$ be the statement that
\[
 (a+b)^n=\sum_{k=0}^{n}\binom nk a^{n-k}b^k
 \quad\text{for every }a,b\in\R.
\]
Part (a) established $P_1$: the two terms in that sum are $a$ and $b$.
Thus this proposition has a valid starting case. The checks at $2$ and
$3$ in (a) are consistent with the result but are not needed as additional
bases for induction.

\step{(c) State the induction hypothesis and the complete next target.}
Let $n\geq1$ be arbitrary and assume $P_n$. Fix arbitrary real numbers
$a$ and $b$. The hypothesis gives
\[
 (a+b)^n=\sum_{k=0}^{n}\binom nk a^{n-k}b^k
\]
for these particular numbers. We must prove
\[
 (a+b)^{n+1}
 =\sum_{k=0}^{n+1}\binom{n+1}{k}a^{n+1-k}b^k.
\]
Because $a$ and $b$ were chosen arbitrarily, obtaining this equality
will prove the assertion for every pair of real numbers, as $P_{n+1}$
requires.

\step{(c) Multiply the assumed identity and distribute both sums.}
Multiply the induction hypothesis by $a+b$. Distributivity of finite
sums gives
\begin{align*}
 (a+b)^{n+1}
 &=(a+b)(a+b)^n\\
 &=(a+b)\sum_{k=0}^{n}\binom nk a^{n-k}b^k\\
 &=a\sum_{k=0}^{n}\binom nk a^{n-k}b^k
   +b\sum_{k=0}^{n}\binom nk a^{n-k}b^k\\
 &=\sum_{k=0}^{n}\binom nk a^{n+1-k}b^k
   +\sum_{k=0}^{n}\binom nk a^{n-k}b^{k+1}.
\end{align*}
In the first sum, multiplying $a^{n-k}$ by $a$ raises its exponent
to $n-k+1=n+1-k$. In the second sum, multiplying $b^k$ by $b$
raises its exponent to $k+1$. All of these are products of nonnegative
integer powers, so these rules are valid at $a=0$ and $b=0$ as well.

\step{(c) Reindex the second sum with its new bounds.}
In that second sum, set $j=k+1$. Then $k=j-1$, and the bounds change
as follows:
\[
 k=0\ \Longrightarrow\ j=1,
 \qquad
 k=n\ \Longrightarrow\ j=n+1.
\]
Also the power of $a$ becomes
\[
 n-k=n-(j-1)=n+1-j.
\]
Every old index corresponds to exactly one new index and conversely.
Therefore
\[
 \sum_{k=0}^{n}\binom nk a^{n-k}b^{k+1}
 =\sum_{j=1}^{n+1}\binom n{j-1}a^{n+1-j}b^j.
\]
Substituting the reindexed sum into the preceding expression gives
\[
 (a+b)^{n+1}
 =\sum_{k=0}^{n}\binom nk a^{n+1-k}b^k
  +\sum_{j=1}^{n+1}\binom n{j-1}a^{n+1-j}b^j.
\]

\step{(c) Separate the two endpoints before combining interior terms.}
The first sum contains an endpoint at $k=0$, equal to
\[
 \binom n0 a^{n+1}b^0=1\cdot a^{n+1}\cdot1=a^{n+1}.
\]
The second sum contains an endpoint at $j=n+1$, equal to
\[
 \binom nn a^0b^{n+1}=1\cdot1\cdot b^{n+1}=b^{n+1}.
\]
Removing these endpoints leaves the common index range $1,\ldots,n$.
Rename the dummy index $j$ in the second interior sum as $k$; this
changes only the name of the summation variable. We obtain
\begin{align*}
 (a+b)^{n+1}
 &=a^{n+1}
  +\sum_{k=1}^{n}\binom nk a^{n+1-k}b^k
  +\sum_{k=1}^{n}\binom n{k-1}a^{n+1-k}b^k
  +b^{n+1}\\
 &=a^{n+1}
  +\sum_{k=1}^{n}
     \left(\binom nk+\binom n{k-1}\right)a^{n+1-k}b^k
  +b^{n+1}.
\end{align*}
For each interior index $k$, both terms contain the identical monomial
$a^{n+1-k}b^k$; the distributive law combines their coefficients.

\step{(c) Apply part (b) and restore the full target sum.}
Every index in the remaining sum satisfies $1\leq k\leq n$, so
part (b) applies to every coefficient. Thus
\[
 (a+b)^{n+1}
 =a^{n+1}
  +\sum_{k=1}^{n}\binom{n+1}{k}a^{n+1-k}b^k
  +b^{n+1}.
\]
The endpoint computations at the beginning of the exercise, applied
with $n+1$ in place of $n$, give
\[
 \binom{n+1}{0}=1,
 \qquad
 \binom{n+1}{n+1}=1.
\]
Consequently the first and last terms can be written as
\[
 a^{n+1}=\binom{n+1}{0}a^{n+1}b^0,
 \qquad
 b^{n+1}=\binom{n+1}{n+1}a^0b^{n+1}.
\]
Including these two endpoint terms with the interior sum produces
\[
 (a+b)^{n+1}
 =\sum_{k=0}^{n+1}\binom{n+1}{k}a^{n+1-k}b^k,
\]
which is exactly the required next statement for the arbitrary real
numbers $a,b$. Therefore $P_n$ implies $P_{n+1}$.

\answer{Part (a) verifies the first three cases, part (b) establishes
Pascal's identity directly from factorials, and part (c) proves the
binomial theorem for every positive integer $n$ and every $a,b\in\R$
by mathematical induction. The proof includes $a=0$, $b=0$, and
$a=b=0$, since it uses no division by either variable and interprets
zero exponents as constant monomials.}
